\section{Background and Related Work}
\label{sec:related}

Five strands of prior work bear on how a researcher should supervise a team of agents.
Each subsection summarizes one strand and states how this paper builds on it or departs
from it. Table~\ref{tab:related} summarizes the comparison.

\subsection{LLM agents and multi-agent frameworks}
\label{sec:rw-agents}

An LLM agent couples a language model with tools and a loop: it decides on an action,
calls a tool, observes the result and decides again~\cite{yao2023react}. Surveys
organize the design space of single agents, for example into profile, memory, planning
and action modules~\cite{wang2024survey}, and of societies of cooperating
agents~\cite{xi2025rise}. Multi-agent frameworks compose several such
agents. AutoGen builds applications from conversable agents that exchange
messages~\cite{wu2023autogen}; MetaGPT encodes standardized operating procedures and
assigns software-company roles to agents~\cite{hong2023metagpt}; ChatDev organizes
role-playing agents into chat-based development phases~\cite{qian2024chatdev}; CAMEL
studies cooperation between role-playing agents~\cite{li2023camel}; and HuggingGPT uses
one model as a controller that plans tasks and dispatches them to specialist
models~\cite{shen2023hugginggpt}. Separating planning from execution also helps at the
level of a single prompt~\cite{wang2023planandsolve}. Practitioner reports distinguish
\textit{workflows}, whose steps are fixed in code, from \textit{agents}, which direct
their own tool use, and recommend an orchestrator--worker pattern for tasks that cannot
be decomposed in advance~\cite{schluntz2024building}; a production research system built
this way used roughly fifteen times the tokens of a chat interaction~\cite{hadfield2025multiagent}.
Agent behavior is increasingly configured through plain files---project instruction files
such as \texttt{AGENTS.md}~\cite{agentsmd2025} and packaged \textit{skills} with a
\texttt{SKILL.md} entry point~\cite{agentskills2025spec}---and coding agents are evaluated
on real repository tasks~\cite{jimenez2023swebench}. The failures of multi-agent systems
cluster into system-design issues, inter-agent misalignment and inadequate task
verification~\cite{cemri2025multiagent}.

This paper adopts the planner--executor and orchestrator--worker patterns and the use of
plain files as configuration. It departs from the frameworks above in two respects. The
orchestrator is the human researcher, whose review time is treated as the scarce
resource, and work passes between agents as versioned files with an explicit status,
not as conversation, so that every handoff can be read, diffed and accepted or rejected.

\subsection{AI systems that automate science}
\label{sec:rw-science}

A parallel line of work aims to automate research itself. The AI Scientist generates
ideas, runs experiments, writes papers and reviews them automatically~\cite{lu2024aiscientist}.
Coscientist designs, plans and executes chemistry experiments through laboratory
automation~\cite{boiko2023autonomous}, and ChemCrow augments a model with expert chemistry
tools~\cite{bran2024chemcrow}. Agent Laboratory takes a human-provided idea through
literature review, experimentation and report writing~\cite{schmidgall2025agentlab}, and
Co-Scientist uses a multi-agent system to generate and refine
hypotheses~\cite{gottweis2025coscientist}. Critics warn that such tools can produce
illusions of understanding, in which scientists believe they understand more than they
do, and can narrow the questions a field asks~\cite{messeri2024illusions}. Others argue
that risk to people grows with an agent's autonomy and that fully autonomous agents should
not be built~\cite{mitchell2025fully}.

Our framework sits at the opposite end of the autonomy range. It keeps the framing of the
question, the choice between approaches and the acceptance of every reported claim with
the researcher, consistent with the accountability required by venues and
funders~\cite{arxiv2025genai,thorp2023chatgpt,nature2023tools,ec2025living}. It also
starts from a constraint specific to engineering: much of the evidence is a physical
measurement, which an agent can analyze but must never be allowed to stand in for.

\subsection{LLMs in hardware design}
\label{sec:rw-eda}

In electrical and computer engineering, work on LLMs has concentrated on design
artifacts. Benchmarks measure the quality of generated Verilog~\cite{thakur2023benchmarking,liu2023verilogeval};
conversational design studies show an engineer and a model co-designing hardware
through dialogue~\cite{blocklove2023chipchat}; and domain-adapted models support
engineering assistance, script generation and bug summarization in chip
design~\cite{liu2023chipnemo}. These studies ask how capable a model is on a design task.
This paper asks a complementary question: given capable models, how should the
researcher who uses them divide, supervise and verify the work of several at once.

\subsection{Levels of automation and human supervisory control}
\label{sec:rw-loa}

Sheridan and Verplank's ten-level scale for supervisory control of teleoperators is the
origin of most \textit{levels of automation}~\cite{sheridan1978human}. Parasuraman,
Sheridan and Wickens separated the type of function automated (information acquisition,
analysis, decision and action) from the level of automation, and proposed that the level
be chosen by evaluating human-performance consequences together with automation
reliability and the costs of decision and action outcomes~\cite{parasuraman2000model}.
Mixed-initiative interfaces weigh the expected value of acting autonomously against that
of asking the user~\cite{horvitz1999principles}; design guidelines for human--AI
interaction~\cite{amershi2019guidelines} and human-centered AI~\cite{shneiderman2020hcai}
argue that high automation and high human control can coexist. For AI agents, Feng,
McDonald and Zhang define five levels by the role the user plays---operator,
collaborator, consultant, approver and observer---and propose autonomy certificates;
they treat autonomy as a design decision separate from capability~\cite{feng2025levels}.
Visibility measures such as agent identifiers, real-time monitoring and activity logs
make agents' actions inspectable after the fact~\cite{chan2024visibility}. Our own
earlier work described advocate agents that act on a person's behalf using explicit,
machine-readable preferences~\cite{ketter2008advocate}.

A second tradition asks how many semi-autonomous systems one person can supervise.
Crandall et al.\ characterized each robot by its \textit{neglect tolerance}---how long
it can be left unattended before its performance degrades---and by the interaction time
needed to restore it, and used the two to estimate the maximum number of robots one
person can manage, a quantity often called \textit{fan-out}~\cite{crandall2005validating}.
Cummings and Mitchell added the wait times that accrue from cognitive reorientation,
from vehicles queueing for the operator, and from lost situation awareness; including
them lowered predicted operator capacity by up to 67\% in their
simulation~\cite{cummings2008predicting}.

We borrow two ideas from this strand: that the level of automation should follow from the
consequences of the action, and that a supervisor's capacity can be computed from the
attention each subordinate system demands. We depart from it in the object being
supervised. A research agent does not need continuous control; it delivers discrete
artifacts---a reading list, a script, a table, a paragraph---whose cost to the human is
the time needed to verify them. Our assignment rule (Section~\ref{sec:autonomy}) makes the
consequence criterion concrete for research tasks, and our review-load model
(Section~\ref{sec:attention}) is a fan-out model for asynchronous artifacts.

\subsection{Automation bias, LLM failure modes and verification}
\label{sec:rw-verification}

People supervising automation tend toward complacency and automation bias: they monitor
less and accept automated recommendations with less scrutiny, especially when the
automation is usually right~\cite{parasuraman2010complacency,goddard2012automation}. LLM
output adds failure modes of its own: hallucination~\cite{ji2023survey}, fabricated and
erroneous citations~\cite{walters2023fabrication}, sycophancy~\cite{sharma2023sycophancy},
degraded use of information in the middle of long contexts~\cite{liu2024lost},
susceptibility to instructions injected through retrieved content~\cite{greshake2023injection},
and, for coding agents, a propensity to exploit test cases~\cite{zhong2025impossiblebench}.
Using one model to judge another scales review but brings biases of its
own~\cite{zheng2023judging}. Against these, scientific computing has long recommended
version control, scripted analysis, recorded provenance and tests~\cite{sandve2013ten,wilson2014best,wilson2017good},
and computational science distinguishes \textit{verification} (solving the equations
right) from \textit{validation} (solving the right equations)~\cite{oberkampf2010vv}.

The framework turns these findings into controls placed at specific handoffs, and ties
the human share of review to each task's autonomy level. Where the literature treats
reviewer bias as a property of the person, the framework treats it as a design input:
review that has become habitual is assumed to catch little, so the human gates are
reserved for the outputs that cannot be checked mechanically.

\begin{table}[htbp]
  \centering
  \caption{How this paper relates to the main strands of prior work. ``Capacity'' asks
  whether the approach models how much a single human can supervise.}
  \label{tab:related}
  \small
  \begin{tabular}{L{3.0cm}L{2.5cm}L{2.7cm}L{2.5cm}L{2.6cm}}
    \toprule
    Approach & Unit supervised & Human role & Autonomy set by & Capacity \\
    \midrule
    Multi-agent LLM frameworks~\cite{wu2023autogen,hong2023metagpt,qian2024chatdev} &
      Agent conversation & Optional participant or approver & System configuration &
      Not a design variable \\
    Automated-science systems~\cite{lu2024aiscientist,schmidgall2025agentlab,gottweis2025coscientist} &
      Research pipeline & Supplies idea; reads output & High by design &
      Not a design variable \\
    Supervisory control~\cite{sheridan1978human,parasuraman2000model,crandall2005validating} &
      Vehicle or robot & Supervisor & Level scale; consequences & Fan-out from neglect and interaction time \\
    Agent autonomy levels~\cite{feng2025levels} &
      Agent & One of five user roles & Design decision; certificates & Not modelled \\
    \textbf{This paper} &
      Research task and its artifact & Orchestrator; sole acceptor of claims &
      Reversibility, reach, verifiability & Daily review-attention budget \\
    \bottomrule
  \end{tabular}
\end{table}
